% Build in this directory: pdflatex -interaction=nonstopmode -halt-on-error class3_handout.tex
\documentclass[a4paper,10pt,pdflatex,ja=standard,jafont=ipaex-type1,
  magstyle=real,scale=1]{bxjsarticle}
\usepackage{lmodern}
\usepackage{amsmath,amssymb,mathtools}
\usepackage{geometry}
\geometry{reset,left=22mm,right=22mm,top=20mm,bottom=17mm,
  headheight=13pt,headsep=6mm,footskip=10mm}
\usepackage{xcolor}
\usepackage{enumitem}
\usepackage{array,booktabs,tabularx}
\usepackage{fancyhdr}
\usepackage{lastpage}

\definecolor{ink}{HTML}{193A5C}
\definecolor{accent}{HTML}{237D82}
\definecolor{soft}{HTML}{EAF3F3}
\definecolor{muted}{HTML}{52616B}
\newcommand{\coursename}{解析学 II}
\newcommand{\handoutnumber}{3}
\pagestyle{fancy}
\fancyhf{}
\fancyhead[L]{\small\color{ink}\coursename{}・授業要約と演習}
\fancyhead[R]{\small\color{ink}第\handoutnumber 回}
\fancyfoot[C]{\small\color{ink}\thepage/\pageref{LastPage}}
\renewcommand{\headrulewidth}{0.3pt}
\setlength{\parindent}{0pt}
\setlength{\parskip}{2pt}
\setlength{\abovedisplayskip}{5pt}
\setlength{\belowdisplayskip}{5pt}
\setlength{\abovedisplayshortskip}{3pt}
\setlength{\belowdisplayshortskip}{4pt}
\setlist[itemize]{leftmargin=1.8em,itemsep=1pt,topsep=2pt,parsep=0pt}

\newcommand{\handouttitle}[1]{%
  \begin{center}
    {\Large\color{ink}#1}\par\vspace{3mm}
    {\color{accent}\rule{0.88\linewidth}{0.7pt}}
  \end{center}\vspace{-2mm}}
\newcommand{\handout}[2]{%
  \renewcommand{\handoutnumber}{#1}%
  \handouttitle{第#1回\quad #2}}
\newcommand{\topic}[1]{\par\vspace{5pt}%
  {\large\sffamily\mdseries\color{accent}#1}\par\vspace{2pt}}
\newcommand{\keybox}[1]{\par\vspace{3pt}\begingroup
  \setlength{\fboxsep}{4pt}%
  \colorbox{soft}{\parbox{\dimexpr\linewidth-2\fboxsep\relax}{\small #1}}%
  \endgroup\par\vspace{2pt}}
\newenvironment{problems}
  {\begin{enumerate}[leftmargin=2em,itemsep=3pt,topsep=3pt,parsep=0pt,
    font=\color{accent}]}
  {\end{enumerate}}


\newcommand{\examplelabel}[1]{{\sffamily\mdseries\color{ink}#1}}
\newcommand{\problemspace}{\par\vspace{5mm}}
\newcommand{\exercises}{%
  {\color{accent}\hrule height 0.7pt}\vspace{3pt}%
  \topic{演習問題（6題）}%
  {\small 計算過程を記し、代表点の選び方や積分領域も明示しよう。}\par}

\begin{document}
\fontsize{10pt}{13.5pt}\selectfont
\setlength{\abovedisplayskip}{4pt}
\setlength{\belowdisplayskip}{4pt}
\handout{3}{重積分の定義}
\keybox{到達目標：リーマン和と重積分の定義を説明し、上和・下和から可積分性を判断できる。
領域上の積分を、長方形上の0延長の積分として理解する。}

\topic{1変数のリーマン和から2変数へ}
$K=[a,b]\times[c,d]$（$a<b$、$c<d$）を長方形とし、
$N,M$ を正の整数とする。各辺を $N,M$ 等分して
$$
 \Delta x=\frac{b-a}{N},\quad \Delta y=\frac{d-c}{M},\qquad
 x_i=a+i\Delta x,\quad y_j=c+j\Delta y
$$
とおく。小長方形 $Q_{ij}=[x_{i-1},x_i]\times[y_{j-1},y_j]$ の
代表点 $\xi_{ij}=(x_{ij}^*,y_{ij}^*)$ を選び、
$$
 S_{N,M}(f)=\sum_{i=1}^N\sum_{j=1}^M
 f(\xi_{ij})\,\Delta x\,\Delta y
$$
をリーマン和という。高さ $f(\xi_{ij})$ に底面積 $\Delta x\Delta y$ を掛けて足す。
$f\ge0$ なら柱の体積の和になり、負の値は符号を付けて加える。
定数関数 $f\equiv h$ では、代表点によらず $S_{N,M}(f)=h(b-a)(d-c)$ となる。

\topic{重積分と可積分性}
$f:K\to\mathbb R$ を有界とする。$K$ を内部が重ならない有限個の小長方形 $Q$ に分けたものを
分割 $\mathcal P$ と呼び、各 $Q$ の代表点 $\xi_Q\in Q$ に対して
$$
 S_{\mathcal P}(f)=\sum_{Q\in\mathcal P}f(\xi_Q)\operatorname{Area}(Q),\qquad
 |\mathcal P|=\max_{Q\in\mathcal P}\operatorname{diam}(Q)
$$
と定める。ここで $\operatorname{Area}$ は面積、$\operatorname{diam}$ は直径を表す。
$|\mathcal P|\to0$ のとき、分割・代表点によらず同じ有限値 $I$ に収束するなら
$$
 \iint_K f(x,y)\,dx\,dy=I:=\lim_{|\mathcal P|\to0}S_{\mathcal P}(f)
$$
と定義し、$f$ はリーマン可積分であるという。
つまり、任意の $\varepsilon>0$ に対してある $\delta>0$ があり、
$|\mathcal P|<\delta$ ならどの代表点でも $|S_{\mathcal P}(f)-I|<\varepsilon$ となる。
一つの分割・代表点の選び方で極限が存在するだけでは十分ではない。

\topic{上和・下和でリーマン和をはさむ}
有界な $f$ に対し、$m_Q=\inf_{p\in Q}f(p)$、$M_Q=\sup_{p\in Q}f(p)$ とおくと
$$
 L(\mathcal P)=\sum_Qm_Q\operatorname{Area}(Q)
 \le S_{\mathcal P}(f)\le
 \sum_QM_Q\operatorname{Area}(Q)=U(\mathcal P).
$$
$L$ を下和、$U$ を上和という。分割を細かくすると下和は増え、上和は減る。
任意の二つの分割は共通の細分をもつので、どの下和もどの上和以下となる。
したがって、$U(\mathcal P)-L(\mathcal P)$ をいくらでも小さくできることが可積分性の判定条件である。

\topic{連続関数の積分が存在する理由}
$f$ が $K$ 上連続なら、閉じた有界な長方形上で一様連続である。
各 $Q$ で最小値・最大値が存在し、振動の上限を
$$
 \omega_f(\delta)=\sup\{\,|f(p)-f(q)|:p,q\in K,\ |p-q|\le\delta\,\}
$$
とおくと、$\omega_f(\delta)\to0$（$\delta\to0$）である。よって
$$
 0\le U(\mathcal P)-L(\mathcal P)
 \le\operatorname{Area}(K)\,\omega_f(|\mathcal P|)\longrightarrow0.
$$
上和と下和は同じ値に近づき、その間のリーマン和も代表点によらず同じ値に収束する。

\examplelabel{例：$[0,1]^2$ 上の $f(x,y)=x+y$ の振動。}
$n\times n$ 分割では、各小正方形で $M_Q-m_Q=2/n$、面積は $1/n^2$ だから
$$
 U(\mathcal P)-L(\mathcal P)=n^2\cdot\frac2n\cdot\frac1{n^2}=\frac2n\longrightarrow0.
$$

\newpage
% Reserve the upper half of page 2 for the remaining summary.
\noindent\begin{minipage}[b][\dimexpr0.5\textheight-\parskip\relax][t]{\linewidth}
\topic{一般の領域：長方形上に0延長する}
有界な領域 $D$ を含む長方形 $K$ を選び、$f:D\to\mathbb R$ を有界とする。
$D$ の外側で $f$ を0に延長した関数 $\widetilde f$ が $K$ 上可積分なら
$$
 \iint_Df\,dx\,dy:=\iint_K\widetilde f\,dx\,dy,\qquad
 \widetilde f(p)=\begin{cases}f(p),&p\in D,\\0,&p\in K\setminus D.\end{cases}
$$
と定める。値は $D$ を含む長方形の選び方によらない。
$D$ の境界の面積が0で、$f$ が $\overline D$ 上連続なら、この積分は存在する。
領域の指示関数を $\chi_D$（$D$ 上で1、それ以外で0）とすると
$$
 \iint_D1\,dx\,dy=\iint_K\chi_D\,dx\,dy=\operatorname{Area}(D).
$$
$f\ge0$ の積分はグラフの下の体積を表し、一般には正の部分と負の部分の体積の差になる。

\topic{積分の基本性質}
$f,g$ が $D$ 上可積分、$r,s\in\mathbb R$ なら、線形性と大小関係が成り立つ。
$$
 \iint_D(rf+sg)=r\iint_Df+s\iint_Dg,\qquad
 f\le g\ \Longrightarrow\ \iint_Df\le\iint_Dg.
$$
各リーマン和でこれらの関係が成り立つので、極限でも保たれる。

\examplelabel{例：単位正方形上の $f(x,y)=xy$。}
$n\times n$ 分割の中点を選ぶと、すべての $n\ge1$ に対し
$$
 S_{n,n}=\frac1{n^4}
 \left(\sum_{i=1}^n(i-\tfrac12)\right)^2
 =\frac1{n^4}\left(\frac{n^2}{2}\right)^2=\frac14.
$$
$f$ は連続で可積分だから、この極限が重積分の値である。
\keybox{確認点：重積分は代表点によらないリーマン和の極限である。
一般の領域では、関数の連続性だけでなく境界の条件も確認しよう。}
\end{minipage}\par\nointerlineskip
% The lower half of page 2 is reserved for exercises and writing space.
\noindent\begin{minipage}[b][0.5\textheight][t]{\linewidth}
\exercises
\begin{problems}
\item $K=[0,2]\times[0,3]$ を $N=2,M=3$ 等分する。
小長方形の面積と、定数関数 $f\equiv4$ のリーマン和を求めよ。\problemspace

\item $K=[0,1]^2$、$f(x,y)=x+y$ とする。
$n\times n$ 分割で右上の頂点を代表点に取ったリーマン和を計算し、その極限を求めよ。\problemspace

\item $K=[0,1]^2$、$f(x,y)=xy$ を $2\times2$ に分割したとき、
上和・下和・中点のリーマン和を求めよ。\problemspace

\item $D=\{(x,y):0\le x\le1,\ 0\le y\le x\}$ とする。
$\iint_D1\,dx\,dy$ の値を図形の面積から求め、正方形 $[0,1]^2$ 上の0延長を書け。\problemspace

\item 境界の面積が0の有界領域 $D$ 上で、可積分な $f$ が $m\le f\le M$ を満たすとする
（$m,M\in\mathbb R$）。
$m\operatorname{Area}(D)\le\iint_Df\le M\operatorname{Area}(D)$ を上和・下和から説明せよ。\problemspace

\item $[0,1]^2$ 上で、$x,y$ がともに有理数なら $f(x,y)=1$、それ以外では0とする。
代表点を2通りに選び、この関数がリーマン可積分でないことを示せ。\problemspace
\end{problems}
\end{minipage}

\newpage
\handouttitle{第3回\quad 演習問題の解答}
\begin{problems}
\setlength{\itemsep}{10pt}
\item $\Delta x=2/2=1$、$\Delta y=3/3=1$ なので、小長方形の面積は $1$ である。
小長方形は $2\cdot3=6$ 個あり、どの代表点でも $f=4$ だから
$$
 S_{2,3}(f)=6\cdot4\cdot1=24.
$$

\item 代表点は $(i/n,j/n)$（$1\le i,j\le n$）、小正方形の面積は $1/n^2$ である。
$\sum_{i=1}^ni=n(n+1)/2$ を用いると
$$
 S_{n,n}(f)=\frac1{n^3}\sum_{i=1}^n\sum_{j=1}^n(i+j)
 =\frac{2n}{n^3}\cdot\frac{n(n+1)}2=1+\frac1n\longrightarrow1.
$$
$f$ は連続なので、$\iint_{[0,1]^2}(x+y)\,dx\,dy=1$ である。

\item $xy$ は $x,y\ge0$ では各変数について単調に増える。
各小正方形の左下の頂点で最小、右上の頂点で最大となり、面積は $1/4$ だから
$$
 L=\frac14\left(0+0+0+\frac14\right)=\frac1{16},\qquad
 U=\frac14\left(\frac14+\frac12+\frac12+1\right)=\frac9{16}.
$$
中点の座標は各変数について $1/4$ または $3/4$ なので
$$
 S_{\mathrm{mid}}=\frac14\left(\frac1{16}+\frac3{16}+\frac3{16}+\frac9{16}\right)
 =\frac14.
$$
確かに $L\le S_{\mathrm{mid}}\le U$ となる。

\item $D$ は頂点 $(0,0),(1,0),(1,1)$ をもつ直角三角形であり、底辺・高さはともに $1$。
したがって、$f\equiv1$ の0延長は
$$
 \iint_D1\,dx\,dy=\frac12,\qquad
 \widetilde f(x,y)=
 \begin{cases}
  1,&0\le y\le x\le1,\\
  0,&0\le x<y\le1
 \end{cases}
 \quad((x,y)\in[0,1]^2).
$$
対角線を含む境界の面積は0なので、0延長は可積分である。

\item $D$ を含む長方形 $K$ で、$\chi_D$ を $D$ 上で1、それ以外で0の関数とする。
$\widetilde f$ を $f$ の0延長とすると $m\chi_D\le\widetilde f\le M\chi_D$ が成り立つ。
同じ分割の下和・上和もこの大小関係を保つ。
境界の面積が0なので $\chi_D$ は可積分であり、分割を細かくして極限を取れば
$$
 m\operatorname{Area}(D)=\iint_Km\chi_D\,dx\,dy
 \le\iint_K\widetilde f\,dx\,dy
 =\iint_Df\,dx\,dy
 \le\iint_KM\chi_D\,dx\,dy=M\operatorname{Area}(D).
$$
$m<0$ や $M<0$ でも、0延長に指示関数を使えば同じ議論が成り立つ。

\item 有理数も無理数も実数の任意の開区間に存在する。
各小長方形の内部で両座標が有理数の代表点を選ぶと、すべての値が1だから
$S_{\mathcal P}(f)=\operatorname{Area}([0,1]^2)=1$ となる。
一方、各代表点の $x$ 座標を無理数に選べば、すべての値が0だから $S_{\mathcal P}(f)=0$。
どれほど分割を細かくしても二つの和は一致せず、代表点によらない極限は存在しない。
したがって、この関数はリーマン可積分ではない。
\end{problems}
\end{document}
